\documentclass[10pt]{article}

% basic packages
\usepackage[margin=1in]{geometry}
\usepackage{fix-cm}        % scale one design size, as OT1 does --- keeps the title bold
\usepackage[T1]{fontenc}   % accented French terms kept inline in the prose
\usepackage[pdftex]{graphicx}
\graphicspath{{./src/}}
\usepackage{amsmath,amssymb,amsthm}
\usepackage{custom}
\usepackage{subcaption}
\usepackage{tabularx}
\usepackage{xcolor}

\usepackage{pgfplots}
\pgfplotsset{compat=1.15}
\usetikzlibrary{arrows,calc,positioning,patterns,decorations.pathreplacing}
\usepgfplotslibrary{fillbetween}

\renewcommand\qedsymbol{$\blacksquare$}

% page formatting
\usepackage{fancyhdr}
\pagestyle{fancy}

\usepackage{titlesec}
\titleformat*{\subsubsection}{\large\bfseries\sffamily}

\renewcommand{\sectionmark}[1]{\markright{\textsf{\arabic{section}. #1}}}
\renewcommand{\subsectionmark}[1]{}
\lhead{\textbf{\thepage} \ \ \nouppercase{\rightmark}}
\chead{}
\rhead{}
\lfoot{}
\cfoot{}
\rfoot{}
\setlength{\headheight}{14pt}

\linespread{1.03} % give a little extra room
\setlength{\parindent}{0in} % reduce paragraph indent a bit
\setcounter{secnumdepth}{2} % no numbered subsubsections
\setcounter{tocdepth}{2} % no subsubsections in ToC

\begin{document}

% make title page
\thispagestyle{empty}
\bigskip \
\vspace{0.1cm}

\begin{center}
{\fontsize{30}{30} \selectfont \bf \sffamily ECO 3TC41: Introduction à l'économie contemporaine }\\
\vspace{24pt}
%{\fontsize{14}{14} \selectfont \rmfamily Donald Kong} \\
%\vspace{6pt}
%{\fontsize{12}{12} \selectfont \ttfamily dkkl1137@gmail.com} \\
%\vspace{24pt}
\end{center}

{\parindent0pt \baselineskip=15.5pt
  The module teaches the basic concepts of economics in order to discuss
  contemporary problems such as competition, economic rationality, monopoly and
  oligopolies, living standards and macroeconomic indicators, money and monetary
  policy, international trade, and the economics of the environment.

  \medskip
  Chapters 1 to 6 are the « micro » half: how transactions are organised, the
  competitive market, welfare analysis, the long run, market failures, and economic
  rationality and strategic reasoning. Chapters 7 to 9 are the « macro » half: the
  role of the State, the national-accounts aggregates, public finances and the
  steering of the economy; the conjunctural policies France has run since 1970; and
  the present French situation --- the sustainability of the public debt and how to
  reduce the deficit.

  \medskip
  Titles and technical terms are in French, the language of the course and its exam:
  each term is glossed in English once, where it is defined, definitions keep the
  course's own French, and the commentary is in English.
}

% make table of contents
% \newpage
\microtoc%
\newpage

% main content
\end{document}
